% ======================================================================
\section{Testing and benchmarking}

% ======================================================================
\subsection{Testing}

% ----------------------------------------------------------------------
\begin{frame}[fragile]{One infamous bug story}

  The binary search function of the Java standard library used to compute
  the median index (\lstinline{low<=high}) like:
  \begin{lstlisting}
    int mid = (low+high) / 2;
  \end{lstlisting}

  Bug went unsuspected for 15+ years.
  --- {\em What bug? That's the exact maths!}
  \begin{itemize}
  \item Sure. But \lstinline{(low+high)} will overflow as soon as \lstinline{low+high > (1<<31)-1}.
  \item \lstinline{mid} index then gets corrupted. Search fails. Critical software at risk. Worldwide.
  \item Patch with exact same maths, without overflowing:
    \begin{lstlisting}
      int mid = low + (high-low) / 2;
    \end{lstlisting}
  \end{itemize}

  Despite correct maths, bug appeared when machines got enough memory to trigger it.

  \begin{itemize}
  \item Trigger as much edge cases as possible during testing.
  \end{itemize}

\end{frame}


% ----------------------------------------------------------------------
\begin{frame}[fragile]{Testing rationale}

  Basic test typology
  \begin{itemize}
  \item Unit testing: Interface behavior
  \item Functional testing: Basic algorithm behavior
  \item Random testing: Functional testing on long streaks of random data
  \end{itemize}

  Functional testing with a {\em digital twin}
  \begin{itemize}
  \item Use simplest, yet likely slower, equivalent code for validation.
  \item Searching an array is simplest: can be provably correct, but runs in $O(N)$.
  \item Searching a Robing-Hood hashtable is more tedious
    \begin{itemize}
    \item Test results on the table are compared to those on the simplest code for arrays.
    \end{itemize}
  \end{itemize}
  
  Most reasonable way to trigger unsuspected edge cases: run code on random data
  \begin{itemize}
  \item Control RNG seed for test debugging, then use random seeds for real tests.
  \item Also known as {\em fuzz testing}.
  \end{itemize}

\end{frame}


% ----------------------------------------------------------------------
\begin{frame}[fragile]{How we'll do the testing}

  \begin{enumerate}
  \item Have unit and functional tests in \lstinline{test/1-map.c}

    Test template is in \lstinline{test/map-test.h} .

  \item Since our digital twin would be logarithmic search at best, use integers instead.

    Cheapest (constant-time), exact local ground truth.

  \item Also use (possibly random) integers in the benchmark.

    That's our fuzzing-like strategy, report failed lookups.
  \end{enumerate}

  Your basic workflow:
  \begin{enumerate}
  \item Implement map option \lstinline{ds_backend_foo}
  \item Check at least you did not break the basics:

    \lstinline{apman$ make deep-clean && make test/1-map@check}

  \item Track the existence of more subtle bugs with the benchmark on random integers

  \item Verify robustness with ascending/descending integer sequences
  \end{enumerate}
  
\end{frame}


% ======================================================================
\subsection{Benchmarking}

% ----------------------------------------------------------------------
\begin{frame}[fragile]{Benchmarking maps}

  \begin{itemize}
  \item Map benchmarking metrics (minimal set), per mapping:
    \begin{itemize}
    \item Memory overhead compared to using a flat array for key/value storage, in bytes
    \item Time to insert-absent/lookup-present/lookup-absent/remove-present, in nsecs
    \end{itemize}
  \item {\em How many nanoseconds to find a specific needle in a mega-haystack?}
    \begin{itemize}
    \item Micro-benchmarking unreliable, aggregate stats on multiple runs
    \item Design benchmark template accordingly and report lookup errors:
      \begin{enumerate}
      \item Insert $N$ different integers in empty map
      \item Lookup those $N$ integers (known to be in the map)
      \item Lookup $N$ other integers (known {\em not} to be in the map)
      \item Remove the $N$ integers
      \end{enumerate}
    \item Save raw run stats in CSV format
    \end{itemize}
  \item Benchmark rendering: timing plots + radar plot with Inter-Quartile Range
    \begin{itemize}
    \item Map implementations feature a variety of behaviors.
    \end{itemize}
  \item Have a Python script to run benchmark configs and make the plots
    \begin{itemize}
    \item Forget \lstinline{matlab}, welcome to  \lstinline{maplab}!
    \end{itemize}
  \end{itemize}

\end{frame}


